\chapter{Reference: row and column layout}
\label{ch:refman-rowcol}

\begin{aside}
Flex containers in two flavours. \code{h(...)} (or
\code{row}) lays children horizontally; \code{v(...)} (or
\code{column}) stacks them vertically.
\end{aside}

\section{Working example}

\begin{lstlisting}[language=C++]
#include <maya/maya.hpp>

using namespace maya;
using namespace maya::dsl;

struct RowCol {
    struct Model {};
    struct Quit {};
    using Msg = std::variant<Quit>;

    static Model init() { return {}; }
    static auto update(Model m, Msg)
        -> std::pair<Model, Cmd<Msg>> {
        return {m, Cmd<Msg>::quit()};
    }

    static Element view(const Model&) {
        return v(
            t<"a column of rows"> | Bold,
            blank_,
            h(
                t<"left">  | Fg<255, 200, 100>,
                spacer_,
                t<"right"> | Fg<100, 200, 255>
            ),
            h(
                t<"A">, text(" "),
                t<"B">, text(" "),
                t<"C">
            ),
            blank_,
            t<"q to quit"> | Dim
        ) | pad<1> | border_<Round>;
    }

    static auto subscribe(const Model&) -> Sub<Msg> {
        return key_map<Msg>({ {'q', Quit{}} });
    }
};

int main() { run<RowCol>({.title = "rowcol"}); }
\end{lstlisting}

\section{Notes}

\begin{itemize}[leftmargin=*]
  \item \code{spacer\_} expands to fill remaining space.
  \item \code{blank\_} is a blank line; equivalent to
    \code{text(" ")} in a column.
  \item You can nest \code{h(v(...))} and \code{v(h(...))}
    freely; layout is computed top-down.
  \item Both containers accept any number of children via
    parameter packs.
\end{itemize}

\section{Sizing}

\begin{itemize}[leftmargin=*]
  \item Default: children take their natural size.
  \item \code{flex<N>}: child gets \code{N} units of weight.
  \item \code{fixed\_width<N>}, \code{fixed\_height<N>}:
    hard size.
\end{itemize}

\section{Recap}

\begin{recap}
\begin{itemize}[leftmargin=*]
  \item \code{h} for rows, \code{v} for columns.
  \item \code{spacer\_} pushes; \code{blank\_} spaces.
  \item Nesting is free; sizing is via modifiers.
\end{itemize}
\end{recap}

\section*{Workshop}

\begin{enumerate}[leftmargin=*]
  \item Build a header bar with title left, date right.
  \item Make a three-column dashboard.
  \item Add a footer with hint text and a help shortcut.
\end{enumerate}
